\documentclass[11pt,a4paper]{article}
\usepackage[T1]{fontenc}
\usepackage{lmodern}
\usepackage[margin=25mm]{geometry}
\usepackage{amsmath,amssymb,amsthm}
\usepackage[hidelinks]{hyperref}
\setlength{\parindent}{0pt}
\setlength{\parskip}{6pt}
\newcommand{\R}{\mathbb R}
\newcommand{\capEHZ}{c_{\mathrm{EHZ}}}
\newcommand{\sys}{\operatorname{sys}}
\begin{document}
\begin{center}
  {\Large The rotated-pentagon profile}\par
  \small An analytic proof by endpoint interpolation
\end{center}

Let $P_5\subset\R^2$ be the regular pentagon centred at the origin with
circumradius one, and let $R(\theta)$ denote counterclockwise rotation through
$\theta$. In $\R^4=\R_q^2\times\R_p^2$, put
\[
 K_\theta=P_5\times R(\theta)P_5,
 \qquad
 d(\theta)=\min_{k\in\mathbb Z}|\theta-k\pi/5|.
\]
We prove
\[
 \boxed{\quad
 \capEHZ(K_\theta)=\bigl(1+\cos(\pi/5)\bigr)^2\sec d(\theta),
 \qquad
 \sys(K_\theta)=\frac{5+2\sqrt5}{10}\sec^2 d(\theta).
 \quad}
\]
Here $\sys(K)=\capEHZ(K)^2/(2\operatorname{vol}_4(K))$.
The proof uses the Haim--Kislev quadratic program and the independently
known Haim--Kislev--Ostrover capacity
\begin{equation}\label{eq:hko}
 \capEHZ(K_{-\pi/2})
 =2\cos(\pi/10)\bigl(1+\cos(\pi/5)\bigr).
\end{equation}
The endpoint-interpolation argument below was supplied by the thesis reviewer.

\textbf{The quadratic quantity being compared.}
Fix the orientation of $P_5$ by taking its outward unit facet normals to be
\[
 n_i=\left(\cos\left(\frac\pi2+\frac{2\pi i}{5}\right),
           \sin\left(\frac\pi2+\frac{2\pi i}{5}\right)\right),
 \quad 0\le i<5,
 \qquad h=\cos(\pi/5).
\]
Thus $P_5=\{q:\langle n_i,q\rangle\le h\text{ for all }i\}$.
The ten normalized facet normals of $K_\theta$ are
\[
 a_i(\theta)=h^{-1}(n_i,0),\qquad
 a_{i+5}(\theta)=h^{-1}(0,R(\theta)n_i),\qquad 0\le i<5.
\]
We use $\omega_0((q,p),(q',p'))=q\cdot p'-p\cdot q'$.
A candidate consists of an ordered list
$\sigma=(\sigma_1,\ldots,\sigma_m)$ of distinct facet indices and weights
$\beta=(\beta_1,\ldots,\beta_m)$ satisfying
\begin{equation}\label{eq:feasible}
 \beta_k\ge0,\qquad \sum_{k=1}^m\beta_k=1,
 \qquad \sum_{k=1}^m\beta_k a_{\sigma_k}(\theta)=0.
\end{equation}
Its quadratic value, with this ordering of the facets, is
\begin{equation}\label{eq:q}
 Q_{\sigma,\beta}(\theta)
 =\sum_{1\le j<k\le m}\beta_j\beta_k
       \omega_0\bigl(a_{\sigma_j}(\theta),a_{\sigma_k}(\theta)\bigr).
\end{equation}
Let $Q_{\max}(\theta)$ be the maximum of these values over all feasible
candidates. The quadratic-program formula says
\begin{equation}\label{eq:hk}
 \capEHZ(K_\theta)=\frac{1}{2Q_{\max}(\theta)}.
\end{equation}
Thus the task is to bound \emph{every} candidate from above, and then find
one attaining that bound.

\newpage
\textbf{Keeping the candidate fixed while the pentagon rotates.}
The closure equation in \eqref{eq:feasible} separates into
\[
 \sum_{\sigma_k<5}\beta_k n_{\sigma_k}=0,
 \qquad
 R(\theta)\sum_{\sigma_k\ge5}\beta_k n_{\sigma_k-5}=0.
\]
Since $R(\theta)$ is invertible, the second equation is equivalent to the
same equation with $R(\theta)$ removed. Consequently, the feasible orders
and weights do not depend on $\theta$. A candidate chosen at one angle
remains feasible, with exactly the same weights and order, at every angle.
We can therefore use its values at two other angles to bound its value
at the angle of interest.

For such a fixed candidate, the pairings in \eqref{eq:q} between normals
from the same factor vanish. Each remaining pairing is, up to sign,
$h^{-2}n_i\cdot R(\theta)n_j$. Hence there are real numbers
$A_{\sigma,\beta}$ and $B_{\sigma,\beta}$, independent of $\theta$, such that
\begin{equation}\label{eq:harmonic}
 Q_{\sigma,\beta}(\theta)
 =A_{\sigma,\beta}\cos\theta+B_{\sigma,\beta}\sin\theta.
\end{equation}
This holds for every feasible candidate; the weights need not maximize
anything.

\textbf{The endpoint bounds control the whole interval.}
We first work on $-\pi/10\le\theta\le\pi/10$.
The pentagon's rotational symmetry gives
$K_{-\pi/2}=K_{-\pi/10}$, since the angles differ by $2\pi/5$.
If $S$ is reflection in a symmetry axis of $P_5$, the map
$(q,p)\mapsto(Sq,Sp)$ is symplectic and takes $K_\theta$ to $K_{-\theta}$.
Thus \eqref{eq:hko} and \eqref{eq:hk} give
\[
 Q_{\max}(-\pi/10)=Q_{\max}(\pi/10)
 =\frac{1}{4\cos(\pi/10)\bigl(1+\cos(\pi/5)\bigr)}.
\]
Every fixed feasible candidate is at most this value at both endpoints.
Its harmonic form \eqref{eq:harmonic} gives the exact identity
\begin{align*}
 Q_{\sigma,\beta}(\theta)
 &=\frac{\sin(\pi/10-\theta)}{\sin(\pi/5)}
       Q_{\sigma,\beta}(-\pi/10)\\
 &\hspace{1em}+\frac{\sin(\pi/10+\theta)}{\sin(\pi/5)}
       Q_{\sigma,\beta}(\pi/10).
\end{align*}
Both coefficients are nonnegative on the interval, and their sum is
$\cos\theta/\cos(\pi/10)$. Substituting the two endpoint bounds therefore
preserves the inequality and yields
\[
 Q_{\sigma,\beta}(\theta)
 \le\frac{\cos\theta}
 {4\cos^2(\pi/10)\bigl(1+\cos(\pi/5)\bigr)}
 =\frac{\cos\theta}{2\bigl(1+\cos(\pi/5)\bigr)^2},
\]
where $2\cos^2(\pi/10)=1+\cos(\pi/5)$.
Taking the maximum over candidates gives the same bound for
$Q_{\max}(\theta)$. Through \eqref{eq:hk}, this is the required
\emph{lower} bound on capacity.

\newpage
\textbf{A candidate attaining the bound.}
Consider the order and weights
\[
 \sigma=(3,8,1,0,5,6),\qquad
 \beta=\left(\frac{\sqrt5}{10},\frac{\sqrt5}{10},
 \frac{5-\sqrt5}{20},\frac{5-\sqrt5}{20},
 \frac{5-\sqrt5}{20},\frac{5-\sqrt5}{20}\right).
\]
The weights are positive and sum to one. The elementary identity
\[
 n_0+n_1=-2h n_3,
 \qquad
 2h\,\frac{5-\sqrt5}{20}=\frac{\sqrt5}{10}
\]
proves closure separately in each factor, so this candidate is feasible
at every angle.

Its value can be checked without expanding all fifteen pairings.
Consecutive entries from the same factor have zero mutual pairing;
by bilinearity, they may be grouped when computing \eqref{eq:q}.
Set
\[
 u=\frac{\sqrt5}{10h}n_3,\qquad v=R(\theta)u.
\]
In the chosen order, the four groups of weighted facet normals are
\[
 (u,0),\quad(0,v),\quad(-u,0),\quad(0,-v).
\]
Their ordered symplectic pairings sum to $2u\cdot v$, giving
\[
 Q_{\sigma,\beta}(\theta)
 =2\left(\frac{\sqrt5}{10h}\right)^2\cos\theta
 =\frac{\cos\theta}{2\bigl(1+\cos(\pi/5)\bigr)^2}.
\]
The last equality uses $h=(1+\sqrt5)/4$.
This attains the upper bound already proved for all candidates. We have
therefore established, including the endpoints,
\[
 Q_{\max}(\theta)
 =\frac{\cos\theta}{2\bigl(1+\cos(\pi/5)\bigr)^2},
 \qquad
 \capEHZ(K_\theta)=\bigl(1+\cos(\pi/5)\bigr)^2\sec\theta
 \quad (|\theta|\le\pi/10).
\]

\textbf{All angles and the systolic ratio.}
The capacity is even in $\theta$, as above, and has period $2\pi/5$.
There is also the symplectic map $(q,p)\mapsto(p,-q)$. It sends $K_\theta$
to $R(\theta)P_5\times(-P_5)$; rotating both factors by $-\theta$ gives
$K_{\pi/5-\theta}$, because $-P_5=R(\pi/5)P_5$.
Together with evenness, this gives period $\pi/5$. Every angle can thus
be reduced to $d(\theta)\in[0,\pi/10]$, proving the stated capacity formula.

Finally, rotation preserves area, and
\[
 \operatorname{vol}_4(K_\theta)=\operatorname{area}(P_5)^2,
 \qquad \operatorname{area}(P_5)=\frac52\sin(2\pi/5).
\]
Substitution into the definition of $\sys$, followed by the elementary
pentagon trigonometric identities, gives
\[
 \sys(K_\theta)
 =\frac{\bigl(1+\cos(\pi/5)\bigr)^4}
 {2\bigl(\frac52\sin(2\pi/5)\bigr)^2}\sec^2 d(\theta)
 =\frac{5+2\sqrt5}{10}\sec^2 d(\theta).
 \qquad\qedsymbol
\]
\end{document}
